\documentclass[11pt]{article}

\usepackage[margin=1in]{geometry}
\usepackage{hyperref}
\usepackage{booktabs}
\usepackage{amsmath}
\usepackage{amssymb}
\usepackage{parskip}
\usepackage{microtype}
\usepackage{setspace}
\usepackage{titlesec}
\usepackage{graphicx}
\usepackage{subcaption}
\usepackage{natbib}
\usepackage{algorithm}
\usepackage{algpseudocode}
\usepackage{multirow}

\hypersetup{
  colorlinks=true,
  linkcolor=black,
  urlcolor=blue,
  citecolor=black
}

\title{\textbf{Performance Cost Regression for Firefox-Blocked Tracker Requests\\via Tweedie Loss and Learned URL Embeddings}}
\author{James Han\\Privacy Team, Mozilla}
\date{January -- April 2026}

\begin{document}

\maketitle

\begin{abstract}
Firefox's Enhanced Tracking Protection blocks third-party tracker requests using the Disconnect list, but currently reports only the number of trackers blocked, with no estimate of the performance cost prevented. We investigate whether machine learning can provide accurate per-request cost estimates for blocked resources. Our work proceeds in two phases. We first develop a domain-level scoring pipeline using XGBoost trained on 348 million HTTP requests and 16 million Lighthouse audits from the HTTP Archive. Critical analysis reveals that this formulation does not require a learned model: domains with CPU measurements need no prediction, and domains without measurements are uniformly inexpensive. We then reformulate the problem as \emph{per-request transfer size prediction}, where the model predicts the response size of a blocked request from pre-response features available at block time---the URL, resource type, initiator, and request metadata. This formulation presents a genuine prediction gap: the response is never observed because the request was blocked. On 348,909 requests across 2,143 tracker domains, we evaluate ten model architectures including tree ensembles, a character-level URL CNN, and specialized loss functions. Individual per-request predictions are inherently noisy due to server-side variation (A/B tests, geo-targeting, SDK version changes invisible in pre-response features), but errors are roughly symmetric and cancel in aggregation. When summed over a typical week of browsing (200 blocked requests), the XGBoost model with Tweedie loss produces a total estimate with median error of 7\%, compared to 21\% for the lookup table baseline; the weekly total is within 10\% of truth 63\% of the time versus 15\% for the lookup table. The Tweedie loss outperforms squared error by 15 percentage points on the same features, demonstrating that loss function selection dominates architecture selection for this zero-inflated target distribution. A character-level convolutional neural network achieves the highest ranking accuracy (Spearman $\rho = 0.977$) by learning URL representations directly from character sequences, replacing hand-crafted regex features with 192-dimensional continuous embeddings.
\end{abstract}

\tableofcontents
\newpage

%% ============================================================
\section{Introduction}
\label{sec:introduction}
%% ============================================================

Firefox's Enhanced Tracking Protection (ETP) blocks third-party tracker requests using the Disconnect tracking protection list~\citep{disconnectlist}, which categorizes domains by privacy function: Advertising, Analytics, Social, Fingerprinting, and Cryptomining. The privacy dashboard at \texttt{about:protections} reports the number of trackers blocked per week, broken down by category. However, the current display treats every blocked tracker equally: a 258-byte tracking pixel from \texttt{bat.bing.com} receives the same weight as a 170\,KB JavaScript SDK from \texttt{googletagmanager.com} that consumes 305\,ms of main-thread CPU time.

This distinction matters to users. The Ghostery Tracker Tax study~\citep{ghostery2020tracker} measured the top 500 websites and found pages loaded in approximately 8.6 seconds with trackers blocked versus 19.3 seconds without, with each additional tracker adding roughly 2.5\% to page load time. For users, a concrete message---``Firefox saved you 2.3\,MB of bandwidth and 890\,ms of CPU time this week''---is more informative and motivating than ``47 trackers blocked.''

Producing such estimates requires predicting the cost of resources that were never fetched. When Firefox blocks a request, the response is never observed: the transfer size, content type, and execution time are unknown. The browser knows only what was visible before the block decision---the URL, resource type, initiator, and request metadata. This information asymmetry defines the prediction problem.

We initially approached this as a \emph{domain-level} scoring problem: assign each tracker domain a static performance cost score using features aggregated from HTTP Archive crawl data~\citep{httparchive}. After building and evaluating the complete pipeline, we identified a fundamental issue: the model was solving a problem that did not require machine learning. Domains with Lighthouse CPU measurements needed no prediction, and domains without measurements were uniformly inexpensive. The deliverable---a static JSON lookup table---could be constructed almost entirely from raw measurements.

We then reformulated the problem as \emph{per-request transfer size prediction}. At the request level, there is a genuine prediction gap. The same tracker domain serves resources of vastly different sizes depending on the URL path and resource type: \texttt{googletagmanager.com/gtag/js} returns a 93\,KB script bundle, while \texttt{googletagmanager.com/collect} returns a 0-byte beacon. A domain-level average conflates these, but the URL path---visible at block time---is highly discriminative. The practical lookup table ceiling is approximately 1,500 entries (domain plus resource type); a path-level table would require over one billion bytes and go stale immediately as tracker SDKs update their URL structures.

\subsection{Contributions}

\begin{enumerate}
    \item We identify and document why domain-level tracker cost estimation does not require a learned model, providing diagnostic evidence including feature space mismatch analysis and a corrected interpretation of the self-training null result.
    \item We formulate per-request transfer size prediction as a supervised learning problem with a genuine inference-time information gap, where pre-response features predict post-response quantities for blocked requests.
    \item We establish the lookup table accuracy ceiling through systematic analysis at six granularity levels and demonstrate that XGBoost with Tweedie loss exceeds this ceiling by 39.4\% overall and 57\% on script resources.
    \item We compare ten model architectures---including Ridge regression, Random Forest, XGBoost, LightGBM, CatBoost, and a character-level URL CNN---and show that loss function selection (Tweedie vs.\ squared error) matters more than model architecture.
    \item We validate deployment readiness through calibration analysis, aggregation error simulation, and distribution shift detection, showing the weekly aggregate estimate is accurate enough for user-facing display.
\end{enumerate}


%% ============================================================
\section{Related Work}
\label{sec:related}
%% ============================================================

\subsection{Third-Party Web Performance Measurement}

The HTTP Archive project~\citep{httparchive} crawls millions of URLs monthly using WebPageTest and Chrome Lighthouse, storing per-request HAR files and per-page audit results in Google BigQuery. The annual Web Almanac~\citep{webalmanac2022} reports on third-party prevalence, finding that the median web page loads resources from 21 third-party domains. Hulce's third-party-web project~\citep{thirdpartyweb} maps approximately 2,700 entity names to their associated domains, enabling Chrome DevTools and Lighthouse to attribute performance costs to named entities.

\citet{pourghassemi2021adperf} developed adPerf, which uses browser activity tracing in Chromium to attribute specific CPU costs to individual ad scripts, decomposing ad-related activity into scripting, styling, layout, and painting. Their approach requires instrumented page loads and cannot be applied to blocked requests where no activity occurs.

A key gap in the literature is the prediction of resource cost for requests that are \emph{prevented} from completing. Existing measurement tools require the response to be observed. Our work addresses the complementary problem: estimating cost from pre-response signals alone.

\subsection{Tree-Based Models for Tabular Prediction}

Gradient boosted decision trees remain the dominant approach for structured tabular data at moderate scale. XGBoost~\citep{chen2016xgboost} introduced scalable tree boosting with regularization. LightGBM~\citep{ke2017lightgbm} improved training speed through histogram-based splitting and leaf-wise tree growth. CatBoost~\citep{prokhorenkova2018catboost} provides native handling of categorical features through ordered target statistics. \citet{grinsztajn2022tree} systematically benchmarked tree-based models against deep learning across 45 tabular datasets and found trees dominant below approximately 10,000 samples, with competitive performance at larger scales. Our dataset of millions of requests is well within the regime where gradient boosted trees are expected to perform strongly.

\subsection{Tracker Blocking in Web Browsers}

Firefox Enhanced Tracking Protection~\citep{disconnectlist} blocks requests to domains on the Disconnect tracking protection list. Brave Browser implements Shields with its own curated filter lists. Privacy Badger uses heuristic learning to identify trackers. uBlock Origin applies pattern-matching filter lists including EasyList and EasyPrivacy. All of these tools make binary block/allow decisions. None estimates the performance cost of what was blocked, which is the gap our work addresses.


%% ============================================================
\section{Domain-Level Cost Estimation: Lessons from the First Approach}
\label{sec:domain-level}
%% ============================================================

We initially developed a domain-level regression pipeline to score each tracker domain on two independent performance axes: main-thread CPU cost and network bandwidth cost. This section summarizes the approach, presents the key findings, and explains why the formulation proved unnecessary as a machine learning problem.

\subsection{Problem Setup and Data}

We constructed ground-truth labels from the June 2024 HTTP Archive crawl (mobile client), drawing on four data sources:

\begin{table}[h]
\centering
\small
\begin{tabular}{lll}
\toprule
Source & Data provided & Scale \\
\midrule
\texttt{httparchive.crawl.requests} & Per-request timing and transfer size & 348M rows \\
\texttt{httparchive.crawl.pages} & Per-page Lighthouse audit results & $\sim$16M pages \\
\texttt{httparchive.almanac.third\_parties} & Entity-to-domain mapping & $\sim$2,700 entities \\
Disconnect \texttt{services.json} & Privacy categories & 4,387 domains \\
\bottomrule
\end{tabular}
\caption{Data sources for domain-level cost estimation.}
\label{tab:data-sources}
\end{table}

After aggregation and filtering to domains appearing on at least 10 pages, the dataset contained \textbf{4,592 tracker domains}. Of these, 2,185 (48\%) had Lighthouse CPU data from the \texttt{bootup-time} audit, while 2,407 (52\%) did not. All domains had network transfer data. The Disconnect list matched 3,686 domains (80\%) via suffix matching, which was necessary because the Disconnect list uses base domains (e.g., \texttt{doubleclick.net}) while HTTP Archive records full subdomains (e.g., \texttt{stats.g.doubleclick.net}).

Two independent target scores in $[0, 1]$ were constructed as weighted percentile ranks:
\begin{align}
    \text{cpu\_cost} &= 0.50 \cdot \text{prank}(\text{total\_cpu}) + 0.35 \cdot \text{prank}(\text{scripting}) + 0.15 \cdot \text{prank}(\text{parse\_compile}) \\
    \text{net\_cost} &= 0.50 \cdot \text{prank}(\text{p50\_bytes}) + 0.30 \cdot \text{prank}(\text{bytes\_per\_page}) + 0.20 \cdot \text{prank}(\text{requests\_per\_page})
\end{align}

A planned third target, render-blocking delay, was dropped after data exploration revealed that trackers almost never block rendering. The top 20 render-blocking domains in the crawl were all first-party CSS and CDN resources; the only tracker was a consent manager. This reflects an economic incentive: ad networks optimize for asynchronous loading to preserve viewability metrics~\citep{sullivan2020cwv}.

\subsection{Model and Results}

We trained two independent XGBoost regressors~\citep{chen2016xgboost}, one per target, using 19 request-level features with hyperparameters tuned via Optuna~\citep{akiba2019optuna}. For CPU cost, training used only the 2,185 domains with real Lighthouse data, after discovering that filling missing values with zero degraded performance (honest Spearman $\rho = 0.734$ with fake zeros versus $0.751$ without).

\begin{table}[h]
\centering
\small
\begin{tabular}{llll}
\toprule
Model & CPU cost $\rho$ & Network cost $\rho$ \\
\midrule
Mean prediction & 0.000 & 0.000 \\
Transfer size only & 0.264 & 0.911 \\
Ridge regression & 0.713 & 0.992 \\
\textbf{XGBoost} & \textbf{0.751} & \textbf{0.999} \\
\bottomrule
\end{tabular}
\caption{Baseline comparison for domain-level cost estimation (Spearman $\rho$).}
\label{tab:domain-baselines}
\end{table}

SHAP analysis~\citep{lundberg2017shap} revealed that the size of the largest script a domain serves (\texttt{max\_script\_bytes}, mean $|\text{SHAP}| = 0.056$) was the dominant predictor of CPU cost, while Disconnect privacy categories contributed near-zero importance. Figure~\ref{fig:cpu-vs-network} confirms the independence of the two target axes: advertising domains span the entire cost space, and no single privacy category predicts performance.

\begin{figure}[h]
    \centering
    \includegraphics[width=0.7\textwidth]{output/cpu_vs_network.png}
    \caption{CPU cost versus network cost for 2,185 domains with Lighthouse data, colored by Disconnect privacy category. The scatter demonstrates target independence.}
    \label{fig:cpu-vs-network}
\end{figure}

\subsection{The Fundamental Problem}
\label{sec:fundamental-problem}

Despite reasonable metrics, the domain-level formulation has a structural flaw: \textbf{the prediction problem it solves does not require a learned model.}

\textbf{Domains with CPU data need no prediction.} The 2,185 domains with Lighthouse measurements have directly computed scores. The model exists only to predict scores for the remaining 2,407 domains.

\textbf{Domains without CPU data are uniformly inexpensive.} Lighthouse's \texttt{bootup-time} audit has a reporting threshold of approximately 50\,ms, confirmed by observing that the minimum reported CPU time across all 2,185 labeled domains is exactly 50.0\,ms. The 2,407 unlabeled domains are unlabeled \emph{because} they fall below this threshold. Assigning them a fixed low score is sufficient.

\textbf{The labeled and unlabeled populations are disjoint.} Kolmogorov-Smirnov tests show statistically significant differences ($p < 0.001$) between labeled and unlabeled domains on every feature. The median \texttt{max\_script\_bytes} is 20,338 for labeled domains versus 585 for unlabeled. Sixty-five percent of unlabeled domains serve no scripts at all.

\begin{table}[h]
\centering
\small
\begin{tabular}{lrr}
\toprule
Feature & Labeled median & Unlabeled median \\
\midrule
\texttt{max\_script\_bytes} & 20,338 & 585 \\
\texttt{script\_request\_ratio} & 0.641 & 0.000 \\
\texttt{network\_cost} & 0.607 & 0.372 \\
\texttt{bytes\_per\_page} & 6,281 & 288 \\
\bottomrule
\end{tabular}
\caption{Feature distributions for labeled (CPU data present) versus unlabeled domains.}
\label{tab:feature-mismatch}
\end{table}

\textbf{Self-training failed due to distribution mismatch, not data sufficiency.} The original analysis attributed the self-training null result~\citep{yarowsky1995unsupervised}---five rounds producing zero improvement---to the labeled set being ``already representative.'' The correct diagnosis is that the labeled and unlabeled sets occupy nearly disjoint regions of feature space. The model's minimum quantile interval width on unlabeled domains was 0.427, exceeding the confidence threshold of 0.30, meaning no unlabeled domain received a confident pseudo-label. Furthermore, the model cannot predict scores below its training range: the minimum prediction for unlabeled domains was 0.215, while the correct values should all be below 0.193 (the score corresponding to the 50\,ms Lighthouse floor).

\textbf{The lookup table is the real deliverable.} The 125\,KB JSON lookup table---4,592 domains with point estimates and prediction intervals---could be constructed almost entirely from raw measurements for the labeled half and a fixed low assignment for the unlabeled half, without any learned model.


%% ============================================================
\section{Per-Request Cost Prediction}
\label{sec:per-request}
%% ============================================================

The failure of the domain-level formulation motivates a reformulation at the \emph{request} level, where a genuine prediction gap exists.

\subsection{Problem Formulation}

When Firefox blocks a tracker request, the browser observes a feature vector $\mathbf{x}$ consisting of the request URL, resource type, initiator type, request priority, and other metadata. The response---including the transfer size $y$ in bytes---is never observed because the request was blocked before the response arrived. The prediction task is:
\[
    \hat{y} = f(\mathbf{x}), \quad \mathbf{x} \in \mathcal{X}_{\text{pre-response}}, \quad y \in \mathbb{R}_{\geq 0}
\]
where $f$ is learned from completed (unblocked) requests in the HTTP Archive, for which both $\mathbf{x}$ and $y$ are observed.

This formulation satisfies the criteria for a genuine machine learning problem:

\begin{enumerate}
    \item \textbf{Information gap at inference.} The label (transfer size) is unobservable at inference time because the request was blocked.
    \item \textbf{Features available at inference.} Firefox observes the URL, resource type, initiator, and request metadata before the block decision.
    \item \textbf{Training labels available.} HTTP Archive contains millions of completed tracker requests with full responses.
    \item \textbf{Lookup table infeasible.} The URL path space is unbounded (query parameters, SDK versions, campaign IDs) and changes continuously. A practical lookup table is limited to domain-plus-resource-type granularity.
\end{enumerate}

\textbf{Browser compatibility.} Transfer size is server-determined: the same URL returns the same payload regardless of whether Chrome or Firefox issues the request. Training on HTTP Archive data (collected via Chrome) and deploying in Firefox is therefore valid for predicting transfer size, though not for predicting timing metrics that depend on the rendering engine.

\subsection{Data Collection}

Training data is drawn from the same HTTP Archive June 2024 crawl used in Section~\ref{sec:domain-level}, but at the per-request level rather than aggregated to domains. The source table \texttt{httparchive.crawl.requests} contains 348 million tracker requests across 5,186 third-party domains in the categories: advertising, analytics, social, tag-manager, and consent-provider.

We extracted a 0.1\% stratified sample via hash-based sampling on the concatenation of page URL and request URL, yielding 348,909 requests across 2,143 domains.

\subsection{Feature Engineering}
\label{sec:features}

Features are organized into four groups based on what Firefox observes at block time. Table~\ref{tab:features} presents the complete feature set.

\begin{table}[h]
\centering
\small
\begin{tabular}{llll}
\toprule
Group & Feature & Type & Encoding \\
\midrule
\multirow{3}{*}{Domain} & \texttt{domain\_median\_bytes} & numeric & Target encoded \\
 & \texttt{domain\_type\_median} & numeric & Target encoded \\
 & \texttt{domain\_request\_volume} & numeric & $\log(\text{count})$ \\
\midrule
\multirow{5}{*}{URL structure} & \texttt{path\_depth} & integer & Raw \\
 & \texttt{url\_length} & integer & Raw \\
 & \texttt{num\_query\_params} & integer & Raw \\
 & \texttt{file\_extension} & categorical & One-hot \\
 & \texttt{path\_token\_count} & integer & Raw \\
\midrule
\multirow{5}{*}{URL content} & \texttt{path\_has\_js} & boolean & Binary \\
 & \texttt{path\_has\_collect} & boolean & Binary \\
 & \texttt{path\_has\_image} & boolean & Binary \\
 & \texttt{path\_has\_sync} & boolean & Binary \\
 & \texttt{path\_has\_ad} & boolean & Binary \\
\midrule
\multirow{6}{*}{Request} & \texttt{resource\_type} & categorical & One-hot \\
 & \texttt{initiator\_type} & categorical & One-hot \\
 & \texttt{chrome\_priority} & ordinal & 0--4 \\
 & \texttt{http\_method} & categorical & One-hot \\
 & \texttt{http\_version} & categorical & One-hot \\
 & \texttt{waterfall\_index} & integer & Raw \\
\bottomrule
\end{tabular}
\caption{Features available at block time. All features are computable from information Firefox observes before the response arrives.}
\label{tab:features}
\end{table}

\textbf{Domain target encoding.} The domain name is the most informative single feature---\texttt{googletagmanager.com} scripts are consistently 85--95\,KB while \texttt{mc.yandex.com} scripts average 455 bytes---but with 1,054 levels, it cannot be encoded as a simple integer (an ordinal encoding that implies a nonexistent numerical ordering). We use target encoding: each domain is replaced by its median \texttt{transfer\_bytes} computed from the training set. To prevent leakage during cross-validation, the encoding is recomputed within each fold.

\textbf{URL content signals.} Regular expressions applied to the URL path capture domain-independent patterns. Paths matching \texttt{\.js|sdk|gtm|gtag} are flagged as likely JavaScript bundles; paths matching \texttt{collect|beacon|pixel|track} as likely beacons. These binary features encode structural knowledge about tracker URL conventions.

\subsection{Baseline: Maximum Feasible Lookup Table}
\label{sec:lut-baseline}

Before introducing a model, we establish the accuracy ceiling of progressively more granular lookup tables.

\textbf{Within-domain variance.} The coefficient of variation (CV) of \texttt{transfer\_bytes} within individual domains has a median of 0.50, with the 75th percentile at 3.21 and the 90th at 29.44. Twenty-five percent of domains exhibit more than threefold variation in transfer size across requests. This variance motivates moving beyond domain-level averages.

\textbf{Feature value beyond domain.} Adding each feature independently to a domain-only baseline yields the following MAE reductions on the proof-of-concept test set:

\begin{table}[h]
\centering
\small
\begin{tabular}{lrr}
\toprule
Feature added to domain & MAE reduction & Spearman $\rho$ \\
\midrule
\texttt{url\_path} (exact match) & 80.5\% & 0.943 \\
\texttt{file\_extension} & 23.4\% & 0.817 \\
\texttt{resource\_type} & 19.8\% & 0.929 \\
\texttt{chrome\_priority} & 10.3\% & 0.787 \\
\texttt{initiator\_type} & 6.2\% & 0.765 \\
\bottomrule
\end{tabular}
\caption{Marginal predictive value of each feature beyond domain identity.}
\label{tab:feature-value}
\end{table}

Exact URL path matching provides 80.5\% MAE reduction, but 47.3\% of test-set paths are unseen in training, making it infeasible as a lookup strategy. A model must generalize from URL \emph{structure} rather than memorize exact paths.

\textbf{Lookup table size analysis.} Table~\ref{tab:lut-analysis} compares lookup tables at increasing granularity.

\begin{table}[h]
\centering
\small
\begin{tabular}{lrrrr}
\toprule
LUT granularity & Entries & Size (gzip) & MAE (bytes) \\
\midrule
Global median & 1 & $<$1\,KB & 15,030 \\
Domain & 968 & $\sim$23\,KB & 9,216 \\
Domain + resource type & 1,340 & $\sim$31\,KB & 8,070 \\
Domain + resource type + extension & 843 & $\sim$20\,KB & 7,606 \\
Domain + URL path & 4,997 & $\sim$117\,KB & 6,075 \\
Domain + URL path + type & 5,151 & $\sim$121\,KB & 6,038 \\
\bottomrule
\end{tabular}
\caption{Lookup table accuracy by granularity on the proof-of-concept test set. The domain + URL path table achieves the lowest MAE but extrapolates to $\sim$1.1\,GB on the full dataset with rapidly decaying coverage.}
\label{tab:lut-analysis}
\end{table}

The practical LUT ceiling is \textbf{domain + resource type} ($\sim$1,500 entries, 31\,KB gzipped, MAE = 8,070 bytes). Path-level tables are infeasible at full scale: the estimated 50 million unique (domain, path) combinations would require approximately 1.1\,GB gzipped and would go stale immediately as tracker SDKs update their URL structures.

%% ============================================================
\section{Experimental Setup}
\label{sec:experiment}
%% ============================================================

\subsection{Data Split}

The 348,909 requests are split 70/15/15 into train (244,235 requests), validation (52,337), and test (52,337) by random row sampling. Row-level splitting is appropriate because Firefox will have HTTP Archive statistics for all Disconnect list domains at inference time; the deployment scenario includes domain identity as a known feature. Target encoding for domain features is computed exclusively from the training split. During cross-validation, encoding is recomputed per fold to prevent leakage.

\subsection{Model Architectures}

We compare ten approaches spanning lookup tables, linear models, tree ensembles with different loss functions, and a neural network.

\begin{table}[h]
\centering
\small
\begin{tabular}{llp{5.5cm}}
\toprule
Model & Tuning & Rationale \\
\midrule
Domain + type LUT & None & Baseline: largest practical lookup table \\
Ridge regression & $\alpha$ via CV & Tests linear separability \\
Random forest & 200 trees, depth 15 & Ensemble without boosting \\
XGBoost (squared error) & Optuna, 75 trials & Standard tree boosting \\
LightGBM & Optuna, 75 trials & Leaf-wise tree growth \\
CatBoost & Optuna, 50 trials & Native categorical handling \\
XGBoost (Tweedie) & Optuna, 40 trials & Magnitude-proportional loss \\
XGBoost (two-stage) & Optuna, 80+80 trials & Classify zero/nonzero, then regress \\
XGBoost + URL hashing & Optuna, 50 trials & Hashed URL path tokens \\
Char-level URL CNN & AdamW, 50 epochs & Learned URL representations \\
\bottomrule
\end{tabular}
\caption{Model architectures compared.}
\label{tab:models}
\end{table}

The Tweedie loss belongs to the exponential dispersion family with variance power $p \in (1, 2)$ tuned via Optuna. It penalizes errors proportionally to predicted magnitude, focusing model capacity on the high-cost requests that dominate user-facing aggregate estimates rather than on the mass of near-zero beacons. The two-stage model separates the prediction into a zero/nonzero classifier (F1 = 0.992) and a regressor on nonzero requests. The URL hashing model replaces hand-crafted regex features with a 64-feature HashingVectorizer over path tokens. The character-level CNN (72,705 parameters) processes raw URL characters through three parallel convolutional layers (kernel sizes 3, 5, 7), concatenates the resulting embedding with tabular features, and regresses through an MLP head.


%% ============================================================
\section{Results}
\label{sec:results}
%% ============================================================

\subsection{Model Comparison}

\begin{table}[h]
\centering
\small
\begin{tabular}{lrrrr}
\toprule
Model & Test MAE & vs LUT & Spearman $\rho$ \\
\midrule
Ridge regression & 9,651 & $-$38.8\% & 0.635 \\
Domain + type LUT & 6,951 & --- & 0.925 \\
CatBoost & 5,228 & +24.8\% & 0.694 \\
Char-level URL CNN & 5,320 & +23.4\% & \textbf{0.977} \\
Quantile $p_{50}$ & 5,462 & +21.4\% & 0.892 \\
XGBoost (squared error) & 4,967 & +28.5\% & 0.774 \\
Two-stage (classify + regress) & 4,943 & +28.9\% & 0.899 \\
LightGBM & 4,826 & +30.6\% & 0.790 \\
Random forest & 4,675 & +32.7\% & 0.913 \\
XGBoost + URL hashing & 4,570 & +34.3\% & 0.734 \\
\textbf{XGBoost (Tweedie)} & \textbf{4,216} & \textbf{+39.4\%} & 0.952 \\
\bottomrule
\end{tabular}
\caption{All models ranked by test MAE. The Tweedie model achieves the best MAE; the CNN achieves the best ranking.}
\label{tab:all-results}
\end{table}

The XGBoost model with Tweedie loss achieves the lowest MAE (4,216 bytes, 39.4\% reduction) and the second-best ranking quality ($\rho = 0.952$). Note that overall MAE is dominated by the $\sim$70\% of requests that are near-zero beacons and trivially predictable; the per-resource-type breakdown in Table~\ref{tab:by-type} and the aggregation analysis in Section~\ref{sec:aggregation} provide more informative evaluations. The character-level CNN achieves the highest ranking accuracy ($\rho = 0.977$) but lags on MAE, likely due to limited training data for a 72K-parameter neural network. The loss function is the primary differentiator: XGBoost with Tweedie outperforms XGBoost with squared error by 15 percentage points despite identical architecture and features.

\subsection{Results by Resource Type}

\begin{table}[h]
\centering
\small
\begin{tabular}{lrrrr}
\toprule
Resource type & $n$ & LUT MAE & Best model MAE & Improvement \\
\midrule
Script & 14,440 & 13,131 & 5,668 & \textbf{+56.8\%} \\
Font & 27 & 80,045 & 37,267 & +53.4\% \\
HTML & 9,321 & 2,367 & 1,422 & +39.9\% \\
CSS & 601 & 8,512 & 6,832 & +19.7\% \\
Image & 13,716 & 9,194 & 8,797 & +4.3\% \\
Text & 4,282 & 420 & 409 & +2.6\% \\
Other & 9,769 & 17 & 30 & $-$76.7\% \\
\bottomrule
\end{tabular}
\caption{MAE by resource type. Improvement is concentrated on scripts, fonts, HTML, and CSS.}
\label{tab:by-type}
\end{table}

The model's value is concentrated on high-cost, high-variance resource types. Scripts see 57\% MAE reduction, fonts 53\%, HTML 40\%. For low-variance types (text, other), the LUT's constant-zero prediction is near-optimal and the model provides no benefit.

\subsection{SHAP Analysis}

SHAP (SHapley Additive exPlanations)~\citep{lundberg2017shap} decomposes each prediction into per-feature contributions. We compute SHAP values on 5,000 test requests from the Tweedie model.

\begin{table}[h]
\centering
\small
\begin{tabular}{rlr}
\toprule
Rank & Feature & Mean $|\text{SHAP}|$ \\
\midrule
1 & \texttt{domain\_type\_median} & 3.722 \\
2 & \texttt{domain\_median\_bytes} & 0.477 \\
3 & \texttt{url\_length} & 0.316 \\
4 & \texttt{path\_depth} & 0.235 \\
5 & \texttt{num\_query\_params} & 0.201 \\
6 & \texttt{domain\_volume} & 0.200 \\
7 & \texttt{rt\_other} & 0.171 \\
8 & \texttt{waterfall\_index} & 0.130 \\
9 & \texttt{method\_POST} & 0.121 \\
10 & \texttt{ext\_js} & 0.115 \\
\bottomrule
\end{tabular}
\caption{Top 10 features by mean absolute SHAP value for the Tweedie model.}
\label{tab:shap}
\end{table}

The target-encoded domain statistic (\texttt{domain\_type\_median}) accounts for the majority of model output (mean $|\text{SHAP}| = 3.72$), nearly $8\times$ the next feature. URL structure features (\texttt{url\_length}, \texttt{path\_depth}, \texttt{num\_query\_params}) collectively contribute substantial signal, validating the hypothesis that URL structure carries predictive information beyond domain identity. The \texttt{rt\_other} flag pushes predictions strongly downward when present, correctly identifying near-zero beacon requests. Notably, the problematic features (\texttt{waterfall\_index}, \texttt{priority\_ord}, \texttt{httpv\_*}) rank low, confirming that their removal for Firefox deployment would have bounded impact.

\subsection{Calibration}

To verify that predictions are reliable enough for user-facing display, we bin predictions by size range and compare predicted means to actual means.

\begin{table}[h]
\centering
\small
\begin{tabular}{lrrrr}
\toprule
Prediction range & $n$ & Predicted mean & Actual mean & Within 25\% \\
\midrule
0--100 bytes & 26,462 & 20 & 63 & 51\% \\
1K--5K & 3,250 & 2,459 & 3,153 & 50\% \\
5K--10K & 1,709 & 7,169 & 8,550 & 66\% \\
10K--25K & 4,044 & 17,146 & 18,993 & 59\% \\
50K--100K & 4,777 & 82,147 & 83,234 & 87\% \\
100K--250K & 576 & 145,427 & 171,091 & 75\% \\
\bottomrule
\end{tabular}
\caption{Calibration: predicted versus actual mean by size bin. The model is well-calibrated for script-sized responses (50K--100K bin: 87\% within 25\%).}
\label{tab:calibration}
\end{table}

The model tracks the diagonal across four orders of magnitude, with strongest calibration in the 50K--100K range (mostly JavaScript bundles), where 87\% of predictions fall within 25\% of actual.

\subsection{Aggregation Accuracy}
\label{sec:aggregation}

Users see a weekly aggregate, not individual predictions. We simulate realistic browsing by randomly sampling $N$ blocked requests, summing predictions, and comparing to the true sum across 1,000 trials.

\begin{table}[h]
\centering
\small
\begin{tabular}{lrrrr}
\toprule
Requests/week & Model: med.\ \% err & LUT: med.\ \% err & Model: $<$10\% & LUT: $<$10\% \\
\midrule
50 & 8.9\% & 20.1\% & 52.8\% & 25.9\% \\
100 & 8.4\% & 21.0\% & 57.5\% & 20.1\% \\
200 & 7.1\% & 22.6\% & 62.7\% & 15.0\% \\
500 & 6.2\% & 24.3\% & 69.0\% & 5.3\% \\
\bottomrule
\end{tabular}
\caption{Aggregation accuracy. ``$<$10\%'' is the fraction of simulated weeks where the total estimate is within 10\% of the true total.}
\label{tab:aggregation}
\end{table}

At typical browsing volumes (200 blocked requests per week), the model's weekly total is within 10\% of truth 62.7\% of the time, compared to 15.0\% for the LUT. The gap widens with more requests because the LUT's systematic bias compounds while the model's roughly symmetric errors cancel in aggregation.

\subsection{Prediction Intervals}

Quantile regression provides per-request prediction intervals at the 10th, 50th, and 90th percentiles. The $p_{10}$--$p_{90}$ intervals achieve 85.6\% coverage (exceeding the 80\% nominal target) with median width of 26 bytes for cheap requests and mean width of 15,869 bytes for expensive ones.

\subsection{Distribution Shift}

A binary classifier trained to distinguish training from test features achieves AUC = 0.557, barely above the 0.5 random baseline. This confirms no meaningful distribution shift within the same crawl month. Requests in the highest out-of-distribution quintile show elevated prediction error (MAE 7,275 vs.\ 4,429 for the lowest quintile), suggesting the OOD score has value for flagging unreliable predictions in production.


%% ============================================================
\section{Learned URL Representations (Work in Progress)}
\label{sec:url-cnn}
%% ============================================================

The hand-crafted URL features (regex flags, path depth, URL length) reduce a rich character sequence to a handful of binary and integer features, discarding combinatorial information. A URL like \texttt{/sdk/v3.2.1/bundle.min.js?config=enhanced} contains patterns---version numbers, \texttt{min}, \texttt{bundle}, the combination of these tokens---that no fixed set of regex features can fully capture. We train a character-level convolutional neural network (CNN) that learns URL representations directly from the character sequence.

\subsection{Architecture}

Each character in the URL path is mapped to a 32-dimensional embedding. Three parallel convolutional layers with kernel sizes 3, 5, and 7 process the sequence, capturing patterns at different character scales (trigrams, tokens, compound patterns). Global max pooling produces a fixed 192-dimensional URL embedding regardless of URL length. This embedding is concatenated with the same tabular features used by the tree models (domain target encoding, resource type, initiator) and fed through a two-layer MLP (128 and 64 units) with batch normalization and dropout. The model has 72,705 parameters.

The CNN is trained with AdamW optimizer, cosine annealing schedule, and log-space MSE loss: $\mathcal{L} = \text{MSE}(\log(1 + \hat{y}), \log(1 + y))$, which penalizes relative rather than absolute errors.

\subsection{Results}

\begin{table}[h]
\centering
\small
\begin{tabular}{lrrr}
\toprule
Model & Test MAE & vs LUT & Spearman $\rho$ \\
\midrule
Domain + type LUT & 6,944 & --- & 0.925 \\
XGBoost Tweedie (hand-crafted) & 4,366 & +37.1\% & 0.944 \\
Char-level URL CNN & 5,320 & +23.4\% & \textbf{0.977} \\
\bottomrule
\end{tabular}
\caption{CNN versus tree model comparison.}
\label{tab:cnn-results}
\end{table}

The CNN achieves the highest ranking quality of any model ($\rho = 0.977$), exceeding XGBoost Tweedie (0.944) and the LUT (0.925). The learned character representations capture ordering information that hand-crafted features cannot express: for example, distinguishing \texttt{/sdk/v3.2/tracker.min.js} (large production bundle) from \texttt{/api/v2/collect} (small beacon) through learned character patterns rather than explicit regex rules.

The CNN's MAE (5,320) lags XGBoost Tweedie (4,366) by 21.8\%. Two factors explain this gap. First, the CNN has 72,705 parameters trained on 296,572 examples (roughly 4 examples per parameter), insufficient for a neural network to outperform a well-tuned tree model. Second, the log-space MSE loss optimizes for relative accuracy (good for ranking) rather than the magnitude-proportional penalty of Tweedie loss (good for MAE on skewed distributions).

The CNN's URL embedding produces 192 continuous features, replacing the 20 hand-crafted binary and integer URL features. These continuous embeddings provide a richer representation that tree models could split on at arbitrary thresholds, enabling a hybrid architecture: CNN URL embedding concatenated with tabular features, fed to an XGBoost model with Tweedie loss. Scaling to the full 35 million request dataset (100$\times$ increase) and implementing this hybrid approach is planned as the next phase of this work.


%% ============================================================
\section{Firefox Browser Integration}
\label{sec:integration}
%% ============================================================

\subsection{Features Available at Block Time}

When Firefox's Enhanced Tracking Protection intercepts a request, the \texttt{nsIChannel} and \texttt{nsILoadInfo} interfaces expose the following information before the response arrives:

\begin{itemize}
    \item \textbf{Request URI:} Full URL including path and query parameters, via \texttt{nsIChannel::GetURI()}.
    \item \textbf{Content policy type:} Resource classification (\texttt{TYPE\_SCRIPT}, \texttt{TYPE\_IMAGE}, \texttt{TYPE\_SUBDOCUMENT}, etc.) via \texttt{nsILoadInfo::GetExternalContentPolicyType()}.
    \item \textbf{Initiator:} The principal that triggered the request, via \texttt{nsILoadInfo::GetLoadingPrincipal()}.
    \item \textbf{HTTP method:} GET, POST, etc., via \texttt{nsIHttpChannel::GetRequestMethod()}.
    \item \textbf{Request headers:} Available via \texttt{nsIHttpChannel::VisitRequestHeaders()}.
    \item \textbf{Third-party classification:} Tracker category (ad, analytics, social, fingerprinting) via \texttt{nsIClassifiedChannel}.
    \item \textbf{Waterfall position:} Relative ordering among concurrent requests.
\end{itemize}

These map directly to the feature set in Table~\ref{tab:features}. The content policy type corresponds to \texttt{resource\_type}; the URI provides \texttt{url\_path}, \texttt{file\_extension}, \texttt{path\_depth}, and all URL content signals; the loading principal provides \texttt{initiator\_type}.

\subsection{Model Deployment}

Unlike the domain-level approach, which ships a static JSON lookup table, the per-request model must perform inference at block time. Two deployment strategies are feasible:

\begin{enumerate}
    \item \textbf{ONNX Runtime:} Export the trained tree ensemble to ONNX format and run inference via a lightweight ONNX runtime. Model size for a 200-tree XGBoost ensemble is typically 200--500\,KB.
    \item \textbf{Compiled decision function:} Translate the tree ensemble into a native C++ decision function. This eliminates runtime dependencies at the cost of update flexibility.
\end{enumerate}

Inference latency for a single tree ensemble prediction is on the order of microseconds, well within the tolerance of the request-blocking path. However, since the cost estimate is used for background aggregation (``Firefox saved you X this week'') rather than real-time decision-making, latency constraints are minimal.

Model updates are delivered via Firefox's Remote Settings infrastructure, the same mechanism currently used for the Disconnect list. Monthly retraining on fresh HTTP Archive crawl data ensures the model tracks changes in tracker URL structures and SDK versions.

\subsection{Privacy Metrics Card}

With per-request predictions, the \texttt{about:protections} card can display specific, quantified savings:

\begin{itemize}
    \item ``Firefox blocked 47 trackers this week, preventing approximately 2.3\,MB of network transfers.''
    \item On mobile, emphasize bandwidth: ``Firefox saved 4.1\,MB of mobile data this month.''
    \item Confidence-aware messaging: if most blocked requests have confident predictions, show exact numbers; if many are uncertain, use softer language (``approximately,'' ``at least'').
\end{itemize}

A hybrid deployment may be optimal: use the model for high-variance resource types (scripts, HTML) where it provides significant accuracy gains, and fall back to the domain-plus-type lookup table for low-variance types (beacons, pixels) where the LUT is already near-optimal.


%% ============================================================
\section{Discussion and Limitations}
\label{sec:discussion}
%% ============================================================

\subsection{Methodological Lessons}

\textbf{Loss function selection dominates architecture selection.} XGBoost with Tweedie loss outperforms XGBoost with squared error by 15 percentage points on the same data and features. The choice between XGBoost, LightGBM, Random Forest, and CatBoost produces smaller differences (24.8--32.7\% improvement range). For zero-inflated, skewed targets, the loss function is the primary modeling decision.

\textbf{Verify the prediction gap before building the model.} The domain-level formulation consumed significant engineering effort before we recognized that the problem did not require machine learning. The diagnostic was straightforward: checking whether labels existed for the inference population. If both features and labels are available for all entities at inference time, the problem is data engineering, not machine learning.

\textbf{Value concentration.} The model's improvement is concentrated on script resources (57\% MAE reduction) with negligible improvement on beacons and pixels. A hybrid deployment---model for scripts, lookup table for beacons---captures most of the accuracy gain at lower complexity.

\textbf{Aggregation rescues per-request noise.} Individual predictions have meaningful error (MAE 4,216 bytes), but the weekly aggregate is within 10\% of truth 63\% of the time at typical browsing volumes. This is because per-request errors are roughly symmetric and cancel in summation. The product metric (weekly total) is substantially more reliable than the modeling metric (per-request MAE).

\textbf{Feature encoding matters.} Integer-encoding a 2,143-level categorical feature (domain) caused XGBoost to assign near-zero importance to domain identity. Target encoding---replacing the domain name with its training-set median transfer size---immediately elevated domain features to top-ranking importance. This was detected through SHAP feature importance analysis after unexpectedly weak initial results.

\subsection{Limitations}

\textbf{Transfer size only.} The model predicts transfer size but not CPU execution time. Transfer size is server-determined and browser-independent, making cross-browser training valid. CPU cost depends on the JavaScript engine and device hardware, which differ between Chrome (training) and Firefox (deployment). Extending to CPU prediction would require Firefox-specific telemetry.

\textbf{Temporal stability.} All data comes from a single month (June 2024). Tracker SDKs update URL structures and bundle sizes over time. The distribution shift analysis (AUC = 0.557) shows no shift within the same month, but the degradation rate across months has not been measured. Monthly retraining on fresh HTTP Archive data is the planned mitigation.

\textbf{Chrome metadata mismatch.} Chrome priority levels and initiator classifications may not map exactly to Firefox's internal representations. These features account for approximately 6\% of model importance (verified via feature importance analysis), bounding the impact of any mismatch.

\textbf{Neural network scale.} The character-level CNN achieves the best ranking ($\rho = 0.977$) but lags on MAE, likely due to limited data (348,909 requests for 72,705 parameters). Scaling to the full 35 million request dataset could close this gap.


%% ============================================================
\section{Conclusion}
\label{sec:conclusion}
%% ============================================================

We investigated machine learning approaches to estimating the performance cost of tracker requests blocked by Firefox's Enhanced Tracking Protection. A domain-level scoring approach, while producing reasonable metrics, was found to solve a problem that did not require a learned model---a finding we document as a methodological contribution.

Reformulating the problem at the per-request level reveals a genuine prediction gap: when a request is blocked, its response is unobserved, but pre-response features carry substantial predictive signal. Across ten model architectures evaluated on 348,909 requests, XGBoost with Tweedie loss reduces mean absolute error by 39.4\% over the best practical lookup table (4,216 vs.\ 6,951 bytes), with 57\% improvement on script resources. A character-level URL CNN achieves the best ranking quality ($\rho = 0.977$), demonstrating that learned URL representations capture patterns beyond hand-crafted features. Aggregation simulation confirms the weekly total estimate is accurate enough for user-facing display, with median error of 7\% at typical browsing volumes. The model can be deployed in Firefox as an ONNX export with microsecond inference, enabling the privacy dashboard to show users specific, quantified estimates of the bandwidth savings provided by Enhanced Tracking Protection.


\bibliographystyle{plainnat}
\bibliography{references}

\end{document}
